\documentclass[11pt,a4paper]{article}
\usepackage[T1]{fontenc}
\usepackage[utf8]{inputenc}
\usepackage[margin=22mm]{geometry}
\usepackage{booktabs,tabularx}
\usepackage[hidelinks]{hyperref}
\setlength{\parindent}{0pt}
\setlength{\parskip}{6pt}
\begin{document}
\raggedright
\title{CON\textsuperscript{2}PHYS questionnaire: reproducible analytical answers}
\author{Neuronal analysis}
\date{}
\maketitle
\textbf{Status: analytical report awaiting recording execution.}
The raw recordings and a completed Q01--Q15 result bundle were unavailable in the authoring environment. No numerical estimates, confidence intervals or winning areas are invented. Missing evidence is not evidence of equality.

The executable source is \path{analysis/+con2phys/qNN_*.m}. Run
\texttt{con2phys.run\_questionnaire} in MATLAB with a raw recording to regenerate this report from the script outputs. The accompanying MAT file preserves results and errors. See \texttt{RESAMPLING.md} for schema and configuration.

\textbf{Inference scope.} Default intervals are within-recording conditional intervals: percentile resampling for cluster summaries and stationary resampling for ordered block/trial summaries. Bayesian intervals must be labelled credible intervals. Animal-level claims require paired animal contrasts. Marginal interval overlap does not select a winning area. Bootstrap repetition counts do not measure empirical coverage.


\clearpage

\section*{Q01. Which brain area has the lowest firing rate over the entire recording?}

\textbf{Answer:} Not determined: the corresponding recording analysis has not been executed.

\textbf{Quantitative results:} unavailable; no numerical answer is asserted.

\begin{tabularx}{\linewidth}{XrrX}

\toprule Group & Estimate & Interval & Units \\

\midrule

Area 1 & -- & -- & spikes/s/cluster \\

Area 2 & -- & -- & spikes/s/cluster \\

Area 3 & -- & -- & spikes/s/cluster \\

\bottomrule

\end{tabularx}

\textbf{Analytical approach.} Full-recording count divided by recording exposure for every mapped cluster, including silent clusters; arithmetic mean by area.

\textbf{Interpretation and limitations.} Cluster resampling conditions on this recording. Mapped clusters are not automatically validated single neurons.

\textbf{Source.} \path{analysis/+con2phys/q01_firing_rate.m}

\clearpage

\section*{Q02. Which brain area has the highest broadband LFP power (1-100 Hz) over the entire recording?}

\textbf{Answer:} Not determined: the corresponding recording analysis has not been executed.

\textbf{Quantitative results:} unavailable; no numerical answer is asserted.

\begin{tabularx}{\linewidth}{XrrX}

\toprule Group & Estimate & Interval & Units \\

\midrule

Area 1 & -- & -- & LFP units squared \\

Area 2 & -- & -- & LFP units squared \\

Area 3 & -- & -- & LFP units squared \\

\bottomrule

\end{tabularx}

\textbf{Analytical approach.} Welch power integrated from 1 to 100 Hz in nonoverlapping blocks; equal channel weighting within each block.

\textbf{Interpretation and limitations.} Stationary resampling acts on ordered block-power summaries. Reference choice, omitted tail and channel validity affect the target.

\textbf{Source.} \path{analysis/+con2phys/q02_broadband_power.m}

\clearpage

\section*{Q03. Which brain area (if any) has the highest density of ripples (i.e. "hippocampal" ripples traditionally occurring during sharp wave-ripples)?}

\textbf{Answer:} Not determined: the corresponding recording analysis has not been executed.

\textbf{Quantitative results:} unavailable; no numerical answer is asserted.

\begin{tabularx}{\linewidth}{XrrX}

\toprule Group & Estimate & Interval & Units \\

\midrule

Area 1 & -- & -- & candidate events/min \\

Area 2 & -- & -- & candidate events/min \\

Area 3 & -- & -- & candidate events/min \\

\bottomrule

\end{tabularx}

\textbf{Analytical approach.} 120-200 Hz envelope candidates: robust onset z=2, peak z=5, duration 25-150 ms; union across channels within 20 ms; counts per minute.

\textbf{Interpretation and limitations.} These are candidate ripples. Sharp-wave morphology and artifacts require validation before hippocampal ripple identification.

\textbf{Source.} \path{analysis/+con2phys/q03_ripple_density.m}

\clearpage

\section*{Q04. In which brain area are pairwise spike train interactions strongest at the 100 ms timescale?}

\textbf{Answer:} Not determined: the corresponding recording analysis has not been executed.

\textbf{Quantitative results:} unavailable; no numerical answer is asserted.

\begin{tabularx}{\linewidth}{XrrX}

\toprule Group & Estimate & Interval & Units \\

\midrule

Area 1 & -- & -- & absolute Pearson r \\

Area 2 & -- & -- & absolute Pearson r \\

Area 3 & -- & -- & absolute Pearson r \\

\bottomrule

\end{tabularx}

\textbf{Analytical approach.} 100 ms spike counts; Pearson correlation for each nonconstant cluster pair; mean absolute correlation per time block.

\textbf{Interpretation and limitations.} Shared neurons make pair estimates dependent. Resample block summaries; shared task drive is not synaptic interaction.

\textbf{Source.} \path{analysis/+con2phys/q04_spike_interactions.m}

\clearpage

\section*{Q05. Which brain area pair has the strongest undirected functional connectivity?}

\textbf{Answer:} Not determined: the corresponding recording analysis has not been executed.

\textbf{Quantitative results:} unavailable; no numerical answer is asserted.

\begin{tabularx}{\linewidth}{XrrX}

\toprule Group & Estimate & Interval & Units \\

\midrule

Area 1 \textless{}-\textgreater{} 2 & -- & -- & absolute Pearson r \\

Area 1 \textless{}-\textgreater{} 3 & -- & -- & absolute Pearson r \\

Area 2 \textless{}-\textgreater{} 3 & -- & -- & absolute Pearson r \\

\bottomrule

\end{tabularx}

\textbf{Analytical approach.} 100 ms counts averaged over clusters in each area; absolute correlation between area populations in each time block.

\textbf{Interpretation and limitations.} Compare the three connections with paired estimates. Shared input and unequal cluster counts can affect connectivity.

\textbf{Source.} \path{analysis/+con2phys/q05_undirected_connectivity.m}

\clearpage

\section*{Q06. Which brain area pair has the strongest directed functional connectivity?}

\textbf{Answer:} Not determined: the corresponding recording analysis has not been executed.

\textbf{Quantitative results:} unavailable; no numerical answer is asserted.

\begin{tabularx}{\linewidth}{XrrX}

\toprule Group & Estimate & Interval & Units \\

\midrule

Area 1 -\textgreater{} 2 & -- & -- & log variance ratio \\

Area 3 -\textgreater{} 2 & -- & -- & log variance ratio \\

Area 3 -\textgreater{} 1 & -- & -- & log variance ratio \\

\bottomrule

\end{tabularx}

\textbf{Analytical approach.} 20 ms population counts; detrended, standardized conditional VAR with five lags; log reduced/full residual variance ratio per block.

\textbf{Interpretation and limitations.} Exploratory in-sample predictability, not causal evidence. Check residuals, stationarity and lag choice; the target is mean blockwise predictability.

\textbf{Source.} \path{analysis/+con2phys/q06_directed_connectivity.m}

\clearpage

\section*{Q07. Which brain area has the highest density of putative fast-spiking interneurons based on spike waveform and/or spike train features?}

\textbf{Answer:} Not determined: the corresponding recording analysis has not been executed.

\textbf{Quantitative results:} unavailable; no numerical answer is asserted.

\begin{tabularx}{\linewidth}{XrrX}

\toprule Group & Estimate & Interval & Units \\

\midrule

Area 1 & -- & -- & fraction of classifiable clusters \\

Area 2 & -- & -- & fraction of classifiable clusters \\

Area 3 & -- & -- & fraction of classifiable clusters \\

\bottomrule

\end{tabularx}

\textbf{Analytical approach.} Classify finite negative-trough/following-positive-peak waveforms using configured sample rate and trough-to-peak width threshold; fraction narrow by area.

\textbf{Interpretation and limitations.} Recorded-cluster proportion is not anatomical density or confirmed interneuron identity. Exclude invalid waveforms explicitly.

\textbf{Source.} \path{analysis/+con2phys/q07_fast_spiking.m}

\clearpage

\section*{Q08. In which brain area are spikes most strongly phase-locked to its own LFP in the 4-10 Hz range?}

\textbf{Answer:} Not determined: the corresponding recording analysis has not been executed.

\textbf{Quantitative results:} unavailable; no numerical answer is asserted.

\begin{tabularx}{\linewidth}{XrrX}

\toprule Group & Estimate & Interval & Units \\

\midrule

Area 1 & -- & -- & PPC \\

Area 2 & -- & -- & PPC \\

Area 3 & -- & -- & PPC \\

\bottomrule

\end{tabularx}

\textbf{Analytical approach.} 4-10 Hz zero-phase filtered LFP on a prespecified channel; complex interpolation at spike times; discard two-second edges; PPC for eligible clusters.

\textbf{Interpretation and limitations.} Default minimum is 50 spikes per cluster. Cluster intervals condition on estimated PPC and eligibility; check reference-channel sensitivity.

\textbf{Source.} \path{analysis/+con2phys/q08_phase_locking.m}

\clearpage

\section*{Q09. Which brain area has the highest excitation-inhibition ratio (i.e., the strongest relative excitation)?}

\textbf{Answer:} Not enough information to identify physiological excitation/inhibition from this proxy.

\textbf{Quantitative results:} unavailable; no numerical answer is asserted.

\begin{tabularx}{\linewidth}{XrrX}

\toprule Group & Estimate & Interval & Units \\

\midrule

Area 1 & -- & -- & broad/narrow ratio \\

Area 2 & -- & -- & broad/narrow ratio \\

Area 3 & -- & -- & broad/narrow ratio \\

\bottomrule

\end{tabularx}

\textbf{Analytical approach.} Estimate the narrow-waveform fraction p and transform its interval to broad/narrow count ratio (1-p)/p, reversing endpoints.

\textbf{Interpretation and limitations.} Physiological excitation/inhibition cannot be identified from waveform counts. Boundary samples have no reported empirical ratio interval.

\textbf{Source.} \path{analysis/+con2phys/q09_excitation_inhibition.m}

\clearpage

\section*{Q10. Which brain area has the shortest intrinsic neuronal timescale during baseline activity?}

\textbf{Answer:} Not determined: the corresponding recording analysis has not been executed.

\textbf{Quantitative results:} unavailable; no numerical answer is asserted.

\begin{tabularx}{\linewidth}{XrrX}

\toprule Group & Estimate & Interval & Units \\

\midrule

Area 1 & -- & -- & seconds \\

Area 2 & -- & -- & seconds \\

Area 3 & -- & -- & seconds \\

\bottomrule

\end{tabularx}

\textbf{Analytical approach.} One-second pretrial baseline without trial overlap; across-trial correlations of 50 ms bins; fit A exp(-lag/tau)+C and summarize accepted tau estimates.

\textbf{Interpretation and limitations.} Requires at least 20 baseline trials, six finite lags, tau 25-1000 ms and R-squared at least 0.5. Interval conditions on accepted fits.

\textbf{Source.} \path{analysis/+con2phys/q10_intrinsic_timescale.m}

\clearpage

\section*{Q11. Which brain area contains most information about variable A? Exclude inter-trial intervals. The specific trial segment on which to base the analysis is not given because it cannot be determined a priori for variable A.}

\textbf{Answer:} Not determined: the corresponding recording analysis has not been executed.

\textbf{Quantitative results:} unavailable; no numerical answer is asserted.

\begin{tabularx}{\linewidth}{XrrX}

\toprule Group & Estimate & Interval & Units \\

\midrule

Area 1 & -- & -- & bits \\

Area 2 & -- & -- & bits \\

Area 3 & -- & -- & bits \\

\bottomrule

\end{tabularx}

\textbf{Analytical approach.} Matched fixed cluster subsets; whole-trial mean firing rates; contiguous five-fold ridge decoding of A; mutual information of held-out confusion matrix.

\textbf{Interpretation and limitations.} Joint trial resampling conditions on fitted predictions. This is decoded information, not full neural mutual information. Circular-shift nulls are separate.

\textbf{Source.} \path{analysis/+con2phys/q11_variable_a_information.m}

\clearpage

\section*{Q12. During which trial segment is variable C best decoded?}

\textbf{Answer:} Not determined: the corresponding recording analysis has not been executed.

\textbf{Quantitative results:} unavailable; no numerical answer is asserted.

\begin{tabularx}{\linewidth}{XrrX}

\toprule Group & Estimate & Interval & Units \\

\midrule

Trial start to stimulus & -- & -- & balanced accuracy \\

Stimulus to outcome & -- & -- & balanced accuracy \\

Outcome to trial end & -- & -- & balanced accuracy \\

\bottomrule

\end{tabularx}

\textbf{Analytical approach.} All clusters; mean rates in the three trial segments; common contiguous five-fold decoder splits; balanced accuracy with fixed classes.

\textbf{Interpretation and limitations.} Joint trial resampling preserves segment pairing. Full prediction coverage is required; absent-class replicates fail explicitly. Training uncertainty and duration effects remain separate.

\textbf{Source.} \path{analysis/+con2phys/q12_variable_c_decoding.m}

\clearpage

\section*{Q13. In which brain area is the dimensionality of neural activity highest? Use the segments from stimulus presentation to outcome to answer this question.}

\textbf{Answer:} Not determined: the corresponding recording analysis has not been executed.

\textbf{Quantitative results:} unavailable; no numerical answer is asserted.

\begin{tabularx}{\linewidth}{XrrX}

\toprule Group & Estimate & Interval & Units \\

\midrule

Area 1 & -- & -- & effective dimensions \\

Area 2 & -- & -- & effective dimensions \\

Area 3 & -- & -- & effective dimensions \\

\bottomrule

\end{tabularx}

\textbf{Analytical approach.} Fixed matched cluster subset in each area; 100 ms bins within stimulus-to-outcome trials; covariance participation ratio tr(C)\textasciicircum{}2/tr(C\textasciicircum{}2), recomputed on whole resampled trials.

\textbf{Interpretation and limitations.} Point and replicate estimators use the same subset. Longer trials contribute more bins; raw covariance weights high-variance units.

\textbf{Source.} \path{analysis/+con2phys/q13_dimensionality.m}

\clearpage

\section*{Q14. In which brain area is modularity the lowest? Use the entirety of the recording to answer this question.}

\textbf{Answer:} Not determined: the corresponding recording analysis has not been executed.

\textbf{Quantitative results:} unavailable; no numerical answer is asserted.

\begin{tabularx}{\linewidth}{XrrX}

\toprule Group & Estimate & Interval & Units \\

\midrule

Area 1 & -- & -- & weighted modularity Q \\

Area 2 & -- & -- & weighted modularity Q \\

Area 3 & -- & -- & weighted modularity Q \\

\bottomrule

\end{tabularx}

\textbf{Analytical approach.} Fixed matched cluster subset; positive Pearson spike-count network at 100 ms; weighted greedy modularity in each time block.

\textbf{Interpretation and limitations.} Mean blockwise modularity is not one whole-session graph statistic. Negative edges are discarded; inspect node-count and resolution sensitivity.

\textbf{Source.} \path{analysis/+con2phys/q14_modularity.m}

\clearpage

\section*{Q15. In which brain area is the neural signal most complex? Use the segments from stimulus presentation to trial end to answer this question.}

\textbf{Answer:} Not determined: the corresponding recording analysis has not been executed.

\textbf{Quantitative results:} unavailable; no numerical answer is asserted.

\begin{tabularx}{\linewidth}{XrrX}

\toprule Group & Estimate & Interval & Units \\

\midrule

Area 1 & -- & -- & normalized entropy \\

Area 2 & -- & -- & normalized entropy \\

Area 3 & -- & -- & normalized entropy \\

\bottomrule

\end{tabularx}

\textbf{Analytical approach.} Permutation entropy of order 3 with approximately 10 ms delay; linear detrending; exclude tied patterns; mean across channels per stimulus-to-end trial.

\textbf{Interpretation and limitations.} High entropy may reflect noise. Assess trial length, valid pattern count, embedding parameters and spectral sensitivity.

\textbf{Source.} \path{analysis/+con2phys/q15_signal_complexity.m}

\clearpage\section*{Section 2: participant meta-information}

Confidence in each answer, prior methodological experience, and expected participant consensus are personal ratings on the original 1--10 scales. They are left unfilled for the participant; execution success or interval width is not substituted for these ratings.

\end{document}
